\section{Related Work}
\label{sec:related}

\textbf{Dependency-based regression test selection.} Ekstazi tracks dynamic dependencies of tests on files and reruns the tests whose dependencies changed~\cite{gligoric2015ekstazi}. STARTS computes dependencies statically at class level, without running tests~\cite{legunsen2017starts}, and a study of static techniques found class-level selection comparable to Ekstazi in performance, at the risk of being unsafe in some cases~\cite{legunsen2016static}. These techniques reason about structure. They need build integration and per-language analysis, whereas \tool{} works from the text of diffs and tests and offers language-agnostic semantic mapping without bytecode instrumentation.

\textbf{Learned selection at industrial scale.} Predictive test selection at Facebook learns from historical outcomes and calibrates a selection cutoff to targets of over 95\% of individual failures and over 99.9\% of faulty changes, reducing total infrastructure cost by a factor of two relative to dependency-based selection~\cite{machalica2019predictive}. Its distinction between test recall and change recall, and its treatment of flakiness through retries, inform our evaluation. Google's experience shows the scale of continuous testing and the rarity of failures~\cite{memon2017taming}, and an assessment of heuristics there found that algorithms based on recent transition history performed below expectations, with the problem remaining largely open~\cite{leong2019transition}. Our guard is a history rule of this kind, and its strength varies across projects (92\% on ShardingSphere, 68\% on Sling), which is one more reason to read it with the failure-label threat in mind. At Ericsson, an experience report used test execution history to improve selection and prioritization on a large subsystem~\cite{najafi2019improving}. These approaches depend on history that a new project does not have.

\textbf{Test case prioritization in CI.} Yaraghi et al.\ define 150 features across nine groups, and evaluate ranking models on 25 systems at test-file level~\cite{yaraghi2023tcpci}. Our data and ranker baseline come from that work. A ranker orders tests without saying where to stop, so a budget such as our top 50\% has to be picked; the calibrated threshold addresses that choice, subject to the limits in Section~\ref{sec:rq1}.

\textbf{LLMs on commits and CI.} Liu et al.\ evaluate LLM-based generation of regression tests from commits~\cite{liu2026cleverest}. That work creates new tests for a change, whereas we choose among existing ones; the challenge makes the same distinction~\cite{challenge}. Recent exploratory efforts leverage LLMs for commit comprehension or test prioritization, but typically run prompts per test execution or lack cost-benefit bounds against raw CI runner expenses; \tool{} differs in its amortized offline index and calibrated risk boundary.

\textbf{Conformal methods.} We build on conformal risk control~\cite{angelopoulos2024crc}, its non-exchangeable extension~\cite{farinhas2024nonexchangeable}, weighted conformal prediction~\cite{tibshirani2019covariate,barber2023beyond}, and high-probability risk control~\cite{bates2021rcps}. To our knowledge, this work represents the first application of weighted non-exchangeable conformal risk control to continuous regression test selection.
